% This file should be rename with the speaker's last name. (e.g : smith)

% Write the author names in the order you wish them to appear
% and underline the one who will give the talk.

% Only the speaker's email address is required.

\begin{abstract_online}{Surfaces and their Duals}{%
        \underline{Josef B\"{o}hm}$^{1}$}{[nojo.boehm@pgv.at]
        }{%
        $^1$DUG, ACDCA, Austria}
    
Recently I found among my many papers a one-page article about Dual Surfaces, written by Dr. Richard Morris, Liverpool University, which appeared in the Maths\&Stats journal which was published for several years by the University of Birmingham. It must have been before 2000 because I couldn’t find the article in the MSOR Maths\&Stats archives:

https://journals.gre.ac.uk/index.php/msor/issue/archive and http://icse.xyz/mathstore/node/568-2.html

Dr. Morris wrote: \emph{The dual of a surface is the set of planes tangent to the surface}.

The mapping works as follows:

Any plane  \emph{a x + b y+ c z = d} can be interpreted as a point  (\emph{a,b,c,d}) in $\textbf{RP}^2$. Morris performs a projection of these points into $\mathbb{R}^3$ using the map
\begin{equation}
 (\emph{a, b, c, d}) \longrightarrow \left(\frac{a}{c} , \frac{b}{c} , \frac{d}{c}\right)
\end{equation}
Some pictures in a low quality were presented. This was all.
I will show how to perform this mapping using  \textbf{\emph{DERIVE}} and TI-Nspire CAS as well and present various surfaces together with their duals.
Then I will vary the mapping followed by parameter curves and their duals. More questions came up, like “How does the dual of the dual look like?” This would be a nasty and boring calculation done by hand, but using a CAS …
Finally, I don’t raise the problem from the third dimension up to the fourth, but do it the reverse way: 
I make a step down and try to find the dual of a plane curve. Cusps appear, where do they come from?
So, this short article of Dr. Morris kept me busy some time and brought a lot of pleasure and surprise into my mathematical life, which I would like to share with an auditorium.







\end{abstract_online}
